\section{Models and Methods}
\label{sec:methods}

\subsection{Task formulation}
Given one image, or $n \in \{4, 8\}$ consecutive frames ending at the labelled frame, the system decides whether a person is down. A VLM receives the image(s) and a prompt and emits a short answer, mapped to \textit{yes}, \textit{no} or \textit{unparsed} (Section~\ref{sec:parsing}). Because a hard yes/no pins the operating point wherever the model happens to prefer, we also read its probability of ``yes'': the probability mass on ``yes'' tokens at the first generated position, renormalised over ``yes'' and ``no''. This score, $p_\text{yes}$, permits threshold-free metrics and lets us choose operating points (Section~\ref{sec:thresholds}).

\subsection{Models}
Table~\ref{tab:models} lists the models. The LFM2.5-VL family \cite{amini2025lfm2,liquidai2026lfm25vl450m,liquidai2026lfm25vl3b} spans three sizes built alike: a SigLIP2 NaFlex vision encoder \cite{tschannen2025siglip2} (86\,M parameters in the smallest model, 400\,M in the others) feeds a hybrid language model of 350\,M, 1.2\,B or 2.6\,B parameters. Varying size within one family holds the training recipe and preprocessing roughly constant. Qwen3.5 at 2\,B and 4\,B \cite{qwen2026qwen35} serves as a reference from another family. We chose models that run in llama.cpp \cite{gerganov2023llamacpp} on our device and carry research-permitting licences. Liquid AI publishes official GGUF exports; Qwen does not, so for Qwen3.5 we use Unsloth's community conversions. Qwen3.5 emits a reasoning trace by default; we disable it so that every model answers directly. Appendix~\ref{app:revisions} lists exact model revisions. We verified that our llama.cpp build accepts image input for both families.

\begin{table*}[t]
\centering
\caption{Models evaluated. Disk size is Q8\_0 language-model weights plus the F16 vision projector (GGUF). $^\dagger$Community GGUF conversion (Unsloth); Qwen publishes no official GGUF.}
\label{tab:models}
\footnotesize
\begin{tabular}{lrlllr}
\toprule
Model & Params & Language backbone & Vision encoder & Licence & Disk (MiB) \\
\midrule
LFM2.5-VL-450M & 0.45\,B & LFM2.5-350M (hybrid) & SigLIP2 NaFlex, 86\,M & LFM Open License 1.0 & 362 + 180 \\
LFM2.5-VL-1.6B & 1.6\,B & LFM2.5-1.2B (hybrid) & SigLIP2 NaFlex, 400\,M & LFM Open License 1.0 & 1189 + 814 \\
LFM2.5-VL-3B & 3\,B & LFM2.5-2.6B (hybrid) & SigLIP2 NaFlex, 400\,M & LFM Open License 1.0 & 2742 + 814 \\
Qwen3.5-2B$^\dagger$ & 2.27\,B & hybrid linear/full attention, 24 layers & ViT, 24 layers & Apache 2.0 & 1919 + 637 \\
Qwen3.5-4B$^\dagger$ & 4.66\,B & hybrid linear/full attention, 32 layers & ViT, 24 layers & Apache 2.0 & 4275 + 641 \\
\bottomrule
\end{tabular}
\end{table*}

\subsection{Prompts}
We test three prompts (full text in Appendix~\ref{app:prompts}). Prompt~\promptA{} is a plain question: ``Is any person fallen or lying on the floor? Answer yes or no.'' Prompt~\promptB{} adds a definition of \pd{} and enumerates the postures that do not count, from kneeling and crouching to lying on furniture or under equipment. Prompt~\promptBtwo{} gives only a one-sentence criterion: ``A person is `down' only if their torso is lying on the floor or ground. Is any person down? Answer yes or no.'' We added \promptBtwo{} after the pilot study (Section~\ref{sec:pilot}) showed that listing look-alike postures in \promptB{} increased false alarms; its criterion mirrors the pose baseline's rule. For multi-frame inputs, a fixed prefix states that the images are consecutive frames in time order and asks about the last.

\subsection{Runtimes and quantization}
All main-grid measurements use llama.cpp \cite{gerganov2023llamacpp} with GGUF weights at three precisions: 16-bit (F16, or BF16 for Qwen), Q8\_0 (8-bit) and Q4\_K\_M (mostly 4-bit, some tensors at 6-bit). The vision projector stays 16-bit throughout, so only the language model's precision varies. We run \texttt{llama-server} with one slot, a 4096-token context and greedy decoding, via its OpenAI-compatible API. In the pilot we also ran LFM2.5-VL-450M in bf16 with Hugging Face Transformers on the same device as an unquantized cross-check. We do not evaluate ONNX Runtime.

\subsection{Input resolution}
\label{sec:resolution}
We resize each frame so that its longest side measures the native size, 512 or 256 pixels, with Lanczos filtering, never upscaling. Frames are sent as lossless PNG, so resolution is the only image variable. The model's own preprocessor resizes and tiles the image again before encoding, so the number of image tokens depends on both our resizing and the model; we report it for every configuration.

\subsection{Output parsing}
\label{sec:parsing}
We cap generation at 16 tokens. An answer is \textit{strictly} parsed if it begins with ``yes'' or ``no'' (any case, optional punctuation) and \textit{leniently} parsed if exactly one of the two words occurs in its first sentence; anything else is \textit{unparsed} and counted as an error. The gap between strict and lenient parsing gauges how faithfully each model follows the format instruction.

\subsection{Decision thresholds}
\label{sec:thresholds}
We report two operating points. The \textit{default} uses the generated answer, equivalent to thresholding $p_\text{yes}$ at roughly 0.5. The \textit{calibrated} point thresholds $p_\text{yes}$ at the value attaining 95\% recall on the development set, because a safety alarm should seldom miss a real event. The threshold consumes labelled data but trains no parameters; we report both so readers can see how much of each model's error originates in its default bias.

\subsection{Pose-based baseline}
\label{sec:baseline}
To test whether a VLM is needed at all, we compare against a keypoint rule. YOLO11 pose models \cite{maji2022yolopose,jocher2024yolo11} (nano and small) estimate 17 body keypoints and a bounding box per person. A person is flagged as down when (i) the torso axis, from hip midpoint to shoulder midpoint, lies within $\theta$ degrees of horizontal and (ii) the bounding box is at least $r$ times as wide as tall; for clips, both must hold in at least a fraction $q$ of the frames. Similar rules have been applied to fall detection before \cite{chen2020openposefall}. We tune $\theta \in \{15^\circ, 20^\circ, \dots, 60^\circ\}$, $r \in \{0.6, 0.8, \dots, 2.0\}$ and, for clips, $q$ on the development set only, with the VLMs' calibrated criterion: the best development F1 among settings with at least 95\% development recall, or the best F1 overall if none reaches that recall. The rule needs no language model and costs a fraction of a VLM's compute, so the question is whether the VLMs' greater robustness to hard negatives justifies their latency.
